\documentclass[10pt,a4paper]{amsart}
\usepackage[margin=19mm]{geometry}
\usepackage{amsmath,amssymb,mathtools}
\usepackage[scr=rsfs,cal=euler]{mathalfa}
\usepackage{xfrac}
\usepackage[unicode,hidelinks,bookmarks=false]{hyperref}
\newcommand{\ie}{i.e.\ }
\newcommand{\cf}{cf.\ }
\newcommand{\resp}{respectively\ }
\newcommand{\Subsection}{\subsection}
\providecommand{\nobreakdash}{\nobreak\hbox{-}}
\setlength{\emergencystretch}{2em}
\multlinegap0pt
\usepackage{bm}
\usepackage{bbm}
\usepackage{dsfont}
\usepackage{xspace}
\usepackage{cancel}
\usepackage{mathtools}
\usepackage{amsthm}
\makeatletter
\@ifundefined{theorem}{\newtheorem{theorem}{Theorem}}{}
\@ifundefined{proposition}{\newtheorem{proposition}[theorem]{Proposition}}{}
\makeatother
\newcommand{\subsequence}[2]{#1\langle#2\rangle}
\title[Pairs of square-free arithmetic progressions in infinite wor]{Pairs of square-free arithmetic progressions in infinite words}
\author{Thomas Del\'epine, Pascal Ochem, Matthieu Rosenfeld}
\date{Source: arXiv:2605.29853v1 (CC BY 4.0)}
\begin{document}
\maketitle

\noindent\emph{Source and licence.} Pairs of square-free arithmetic progressions in infinite words. Version arXiv:2605.29853v1 (CC BY 4.0), distributed under the Creative Commons Attribution 4.0 International licence, as recorded in the arXiv licence metadata for this exact version and shown on the abstract page https://arxiv.org/abs/2605.29853v1. Full rights notice: licenses/arxiv-2605.29853-CC-BY-4.0.txt.\par\smallskip

\noindent\textit{Editorial scope.} Counted unit: the Proposition whose source label is \texttt{prop:6-compl-sqfree} in the article (source0.tex lines 561--563, source label \texttt{prop:6-compl-sqfree}), delivered together with the article's own complete original proof of it (source0.tex lines 565--585, one \texttt{proof} environment of 21 lines). No mathematics is written, repaired or re-derived: statement and proof are the article's own text line for line. Required context retained uncounted: none. Every definition, display and label of those lines is retained unchanged. Omitted from this package and not redistributed: 1--560, 564--564, 586--861. Nothing inside a retained excerpt is omitted or altered: every retained body is delivered exactly as the article's lines are written, so no omission receipt of any kind accompanies this package. Citations: the retained excerpts carry no citation command at all, so no bibliography entry of the article is transcribed and no reference is redistributed. Reference adaptation: none was needed and none was made; every reference inside the delivered unit resolves inside it, so no unresolved reference or citation is produced. Printed slips kept literal: no printed word, symbol, hypothesis or number of the counted statement or of its proof was altered. Layout and class: the article's own class is \texttt{article} with packages including inputenc, amsmath, amsthm, amssymb, graphicx, fullpage, mathptmx, helvet, xcolor, hyperref, cleveref, cite, url, mathtools; the wrapper is the lane's pinned \texttt{amsart} editorial wrapper at 10pt on A4, which restates the theorem-like environments the retained text needs and adds the article's own macro definitions that the retained text needs (\texttt{\textbackslash subsequence}), with extra wrapper packages (bm, bbm, dsfont, xspace, cancel, mathtools). The delivered numbering is the unit's own. Assets: no retained span contains a figure environment, a graphics-inclusion command, a PDF-page inclusion, a table or a TikZ picture; no image file is copied into this package, the package declares no asset dependency and no image file is redistributed with it. Licence: version arXiv:2605.29853v1 of this article is distributed under the Creative Commons Attribution 4.0 International licence, as recorded in the arXiv licence metadata for this exact version and shown on the abstract page https://arxiv.org/abs/2605.29853v1. A full rights notice is delivered at licenses/arxiv-2605.29853-CC-BY-4.0.txt. Characters: every retained excerpt is pure ASCII, so the delivered engine, XeLaTeX with its default Latin Modern text font, reports no missing character.
\begin{proposition}\label{prop:6-compl-sqfree}
    For every (infinite) ternary square-free word $t$, $\mathcal{R}(t)$ is an (infinite) ternary square-free word.
\end{proposition}

\begin{proof}
    Assume for the sake of contradiction that $\mathcal{R}$ does not preserve square-freeness. So let $t$ be a finite square-free word such that $\mathcal{R}(t)$ contains a square $uu$ of period $\Delta$ for some $\Delta \geq 1$. A simple inspection of $\mathcal{R}$ shows that $t$ must be of length at least $2$ so there exists a non-empty ternary square-free word $t' \in \Sigma_3^+$ and $x \in \Sigma_3$ such that $t = t'x$.
    Assume that $\Delta \ge 10$. Hence, the first occurrence of $u$ contains at some position $d\in\{0,1,2,3,4,5\}$ the beginning of a block $R_{ij}$ for some $i$, $j \in \{0, 1, 2\}$ which is a factor $abc$ with $b = \pi_3(a)$ and $c = \pi_3(b)$.

    Since $\pi_3\circ\mathcal{R} = \mathcal{R}\circ\pi_3$, we can assume without loss of generality that $abc = 012$. As already noted, $\subsequence{\mathcal{R}(t)}{6} = t'$, so since $t$ is square-free, then so is $\subsequence{\mathcal{R}(t)}{6}$. If $\Delta \equiv 0 \pmod6$, then for all $i \in \{0, 1, 2, 3, 4, 5\}$ the subsequence congruent to $i$ modulo $6$ of $uu$ is a square and for one of these $i$ this sequence is a factor of $\subsequence{\mathcal{R}(t)}{6}$ which is absurd since $\subsequence{\mathcal{R}(t)}{6}$ is square-free.
    
    Therefore, we can assume that $\Delta \equiv \alpha\pmod6$ with $0 < \alpha < 6$. 
    Observe that the factor $012$ only appears at position $0$ of $R_{01}$, at position $0$ of $R_{02}$, at position $1$ of $R_{21}$, at position $5$ of $R_{21}R_{10}$, or at position $5$ of $R_{21}R_{12}$ (remember that by definition of $\mathcal{R}$, $R_{12}R_{10}$ or $R_{12}R_{12}$ cannot appear in the image of a word by $\mathcal{R}$). By hypothesis, the factor $012$ at position $d$ of the first period is created by either the prefix of size $3$ of $R_{01}$ or of $R_{02}$. Since $0 < \alpha < 6$, the factor $012$ at position $d$ of the second period must be created by either $R_{21}$, $R_{21}R_{10}$, or $R_{21}R_{12}$ . 
    \begin{itemize}
        \item If $d \geq 4$, then the factor $012$ at position $d$ of the first period must be preceded by $R_{10}$ or $R_{20}$, so $u$ contains either the factor $102012$ or the factor $1021012$ at position $d - 4$. If the factor $012$ at position $d$ of the second period is created by $R_{21}$, then similarly, it must be preceded by $R_{02}2$ or $R_{12}2$, so $u$ must contain either the factor $0212012$ or the factor $2102012$ at position $d - 4$ which is a contradiction. Otherwise, if the factor $012$ at position $d$ of the second period is created by $R_{21}R_{10}$ or $R_{21}R_{12}$, then $u$ must contain either the factor $0121012$ or the factor $2021012$ at position $d - 4$ which is a contradiction.
        \item If $d < 4$, then the factor $012$ at position $d$ of the first period must be followed by either $102$ or $021$ depending on whether $012$ is created by $R_{01}$ or $R_{02}$, so since $\Delta \geq 10$ and $d < 4$, $u$ must contain either the factor $0121021$ or the factor $0120212$ at position $d$. If the factor $012$ at position $d$ of the second period is created by $R_{21}$, then it is followed by either $10R_{10}$ or $10R_{12}$. In both cases, $u$ must contain the factor $0121012$ at position $d$ which is a contradiction. Otherwise, if the factor $012$ at position $d$ of the second period is created by $R_{21}R_{10}$ or $R_{21}R_{12}$, then $u$ must contain either the factor $0120102$ or the factor $0120210$ at position $d$ which is a contradiction.
    \end{itemize}
    Finally, if $\Delta \leq 9$, the square $uu$ in $\mathcal{R}(t)$ must start at position $d \leq 5$ in the image of a square-free word of length $5$ under $\mathcal{R}$. To verify that no such word exists\footnote{It can also easily be done with a simple computer assisted verification.}, it is sufficient to look at the image of a square-free word of length $4$ starting with $01$ or $02$ under $\mathcal{R}$ and observe, for each of the following words, for each period $\Delta \in \{1, 2, 3, 4, 5, 6, 7, 8, 9\}$ and for each position $d \in \{0, 1, 2, 3, 4, 5\}$, that the factor of length $\min(\Delta, 5)$ starting at position $d$ is always different from the factor of the same length starting at position $d + \Delta$:
    \begin{align*}
        &\textbf{0}12102\textbf{1}20102\textbf{0}12021\textbf{2}, \textbf{0}12102\textbf{1}20210\textbf{2}01021\textbf{0},
    \textbf{0}12102\textbf{1}20210\textbf{2}01210\textbf{1},\\
        &\textbf{0}12021\textbf{2}01021\textbf{0}12102\textbf{1},
    \textbf{0}12021\textbf{2}01210\textbf{1}20102\textbf{0},
    \textbf{0}12021\textbf{2}01210\textbf{1}20210\textbf{2}. \qedhere
    \end{align*}
\end{proof}

\end{document}
